% Source PDF: source/普通高考/2014/2014新课标2理(贵州,甘肃,青海,西藏,黑龙江,吉林,海南,宁夏,内蒙古,新疆,云南).pdf, page 1 (question 7).
% PDF embedded-bitmap attempt: pdfimages -list returned no image objects; the flowchart is vector line art.
% Grid evidence: tmp/review_2014_ncp2_science/grid/page_grid100.png and q7_grid50.png;
% comparison source copied to tmp/redraw_flowchart_2014_new_curriculum_paper_2_science/q7_original_final.png
% from q7_orig_fig_final2.png after the ungridded page render.
% Node -> directed-edge list: 开始 -> 输入(x,t) -> M=1,S=3 -> k=1 -> k<=t;
% 是 -> M=(M/k)x -> S=M+S -> k=k+1 -> left loop-back -> k<=t;
% 否 -> 输出 S -> 结束.
% solid segments: all node outlines and directed connectors; dashed segments: none.
% labels: branch labels sit beside the corresponding decision exits; every node is
% locally zero-indented, horizontally/vertically centered, and compactly padded.

\begin{tikzpicture}[
  x=1cm,
  y=1cm,
  >=Stealth,
  line width=0.45pt,
  line cap=round,
  line join=round,
  font=\normalsize,
  flowtext/.style={anchor=center, align=center,
    execute at begin node={\setlength{\parindent}{0pt}}},
  flowbox/.style={draw, minimum height=0.68cm, inner xsep=4pt, inner ysep=2pt,
    align=center, execute at begin node={\setlength{\parindent}{0pt}}},
  terminal/.style={flowbox, rounded corners=0.18cm, minimum width=1.15cm},
  io/.style={flowbox, trapezium, trapezium left angle=72, trapezium right angle=108,
    minimum height=0.72cm},
  decision/.style={flowbox, diamond, aspect=2.85, minimum width=2.85cm,
    minimum height=0.90cm, inner xsep=1pt, inner ysep=1pt},
  branchlabel/.style={flowtext, inner sep=0pt},
]
  % Main vertical sequence.
  \node[terminal] (start) at (0,10.85) {开始};
  \node[io, minimum width=1.85cm] (input) at (0,9.72) {输入 $x,t$};
  \node[flowbox, minimum width=2.65cm] (init) at (0,8.55) {$M=1,S=3$};
  \node[flowbox, minimum width=1.45cm] (kinit) at (0,7.58) {$k=1$};
  \node[decision] (test) at (0,6.20) {$k\leq t$};

  \node[flowbox, minimum width=2.70cm, minimum height=1.00cm] (scale)
    at (-2.20,4.92) {$M=\dfrac{M}{k}x$};
  \node[flowbox, minimum width=2.40cm] (sum) at (-2.20,3.55) {$S=M+S$};
  \node[flowbox, minimum width=2.20cm] (increment) at (-2.20,2.20) {$k=k+1$};

  \node[io, minimum width=1.85cm] (output) at (3.35,4.92) {输出 $S$};
  \node[terminal] (finish) at (3.35,3.70) {结束};

  \draw[->] (start.south) -- (input.north);
  \draw[->] (input.south) -- (init.north);
  \draw[->] (init.south) -- (kinit.north);

  % One explicit shared stem receives the loop return and enters the decision.
  \coordinate (loop-merge) at (0,7.15);
  \draw (kinit.south) -- (loop-merge);
  \draw[->] (loop-merge) -- (test.north);

  % 是 branch: decision -> M update -> sum -> increment.
  \draw[->] (test.west) -- (-2.20,6.20) -- (scale.north);
  \draw[->] (scale.south) -- (sum.north);
  \draw[->] (sum.south) -- (increment.north);

  % 否 branch: independent right exit to output and finish.
  \draw[->] (test.east) -- (3.35,6.20) -- (3.35,5.40) -- (output.north);
  \draw[->] (output.south) -- (finish.north);

  % Loop from k=k+1 rises on the left and returns to the shared stem.
  \draw (increment.west) -- (-3.90,2.20) -- (-3.90,6.95);
  \draw (-3.90,7.15) -- (-0.60,7.15) -- (loop-merge);

  \node[branchlabel, anchor=south] at (-1.52,6.47) {是};
  \node[branchlabel, anchor=south] at (2.52,6.47) {否};
\end{tikzpicture}
